% ЗӨВХӨН БҮТЭН НООРОГТ (main.tex) хэвлэгдэнэ. Багшид үзүүлэх ноорогт (main_teacher.tex) орохгүй.
% Хариулт байхгүй асуулт (дэд хэсэг) ба Мэдлэгийн мод дээрх хайлт
\subsection{Хариулт байхгүй асуулт}
\label{sec:lit-unanswerable}
% SQuAD 2.0 (RAG08): агент бүтэн уншиж papers/16, papers/22-т бүртгэсэн. Хэрэглэгч W03-т уншина.

Сургуулийн туслахад хүн мэдээллийн санд байхгүй зүйлийг ч асууна. Ийм үед систем таамаглаж хариулах нь аюултай. Rajpurkar нар хариулт нь өгөгдсөн бичвэрт байхгүй 53 мянга гаруй асуулттай SQuAD 2.0 өгөгдлийг гаргасан \cite{rajpurkar2018squad2}. Асуултууд тухайн бичвэртэй холбоотой, үнэмшилтэй харагддаг боловч хариулт нь бичвэрт байхгүй. Өмнөх SQuAD 1.1 өгөгдөл дээр F1 (нарийвчлал ба бүрэн байдлыг нэгтгэсэн оноо) 85.8 авдаг загварын бүтэц хариулт байхгүй асуулт нэмэгдсэн SQuAD 2.0 дээр 66.3 авсан бол хүн 89.5 авсан \cite{rajpurkar2018squad2}. Хэрэглэгчийн зорилгыг ангилдаг системүүд ч хамрах хүрээндээ байгаа асуултыг сайн ангилдаг боловч хүрээнээс гадуурх асуултыг танихдаа сул байдаг \cite{larson2019clinc}. Marcel системд зөв баримтыг өгсөн үед хариулах боломжгүй асуултаас зөв татгалзах нь 0.68-аас 0.93 болж сайжирсан \cite{trienes2025marcel}.

Эдгээр ажил хариулахаас татгалзахыг хэмжсэн бөгөөд даалгавар, хэл нь энэ ажлынхаас өөр. Энэ ажилд систем татгалзаад зогсохгүй, хандах алба, өрөөг заах ёстой. Иймд цуглуулах асуултаас хариулт нь сургуулийн мэдээлэлд байхгүй асуултыг тусад нь тэмдэглэж, ийм үед систем хандах алба, өрөөг зөв зааж байгаа эсэхийг шалгана (\ref{sec:exp-tree}-р хэсэг).

\section{Мэдлэгийн мод дээрх хайлт}
\label{sec:lit-tree}

Энэ ажилд мэдлэгийн мод гэж сургуулийн мэдээллийг салаагаар ангилсан бүтцийг хэлнэ. Жишээ нь «Элсэлт $\rightarrow$ Төлбөр», «Байршил $\rightarrow$ Өрөө». Энэ хэсэгт ийм модыг хайлтад хэрхэн ашиглаж болохыг авч үзэв.

Модонд суурилсан хайлтын түгээмэл аргын нэг нь RAPTOR юм \cite{sarthi2024raptor}. Энэ арга бичвэрийн хэсгүүдийг автоматаар бүлэглээд, бүлэг бүрийн хураангуйг хэлний загвараар бичиж, модыг доороос дээш босгодог. Хураангуйг хайлтын нэмэлт зангилаа болгон ашигладаг. Энэ ажилд ийм автомат мод босгоогүй нь хоёр шалтгаантай. Нэгдүгээрт, ресепшний ажилтны нэрлэсэн сэдвүүд (\tref{tab:reception_topics}) төлбөр, өрөө, холбоо барих хүн зэрэг нэг баримтаар хариулагдах асуулт гэж үзсэн; үүнийг цуглуулах асуултаар шалгана. Нэг баримтын асуултад энгийн хайлт дутаагүй: Han нарын үнэлгээнд ийм асуулттай Natural Questions өгөгдөл дээр энгийн RAG F1 64.78, RAPTOR-т суурилсан арга 60.04 авсан \cite{han2025ragvsgraphrag}. Харин олон баримтыг холбох асуултад модтой арга ялимгүй илүү байсан бөгөөд зохиогчид хоёр төрлийн аргыг бие биеэ нөхдөг гэж дүгнэсэн. Хоёрдугаарт, хэлний загварын бичсэн хураангуй буруу мэдээлэл агуулж болдог. RAPTOR-ын зохиогчид гараар шалгасан 150 зангилааны 4\%-д буруу мэдээлэл байсан \cite{sarthi2024raptor}.

Хүний бэлдсэн ангиллыг хайлтад ашигласан ажил ч бий. TaxoIndex эрдэм шинжилгээний өгүүллийн хайлтад бэлэн салбарын ангиллыг индекст ашигласан \cite{kang2024taxoindex}. T-RAG нь нэг байгууллагын бүтцийн модыг асуултад нэмэлт мэдээлэл болгон ашиглаж, 37 асуултын 27-д зөв хариулсан бол ердийн RAG 21-д зөв хариулсан \cite{fatehkia2024trag}. Эдгээрийн үнэлгээний багц жижиг эсвэл өөр салбарынх тул үр дүнг энэ ажилд шууд шилжүүлэх боломжгүй.

Энэ ажилд гараар хийсэн модыг хайлтад ашиглах гурван хувилбарыг авч үзэв (\tref{tab:tree_options}).

\begin{table}[H]
\centering
\caption{Мэдлэгийн модыг хайлтад ашиглах гурван хувилбар}
\label{tab:tree_options}
\begin{tabular}{|p{2.6cm}|p{5.7cm}|p{5.3cm}|}
\hline
\textbf{Хувилбар} & \textbf{Юу хийх} & \textbf{Эрсдэл} \\
\hline
Энгийн хайлт & Бүх эх сурвалжаас хайна. Мод зөвхөн хариулт олдохгүй үед хэнд, хаана хандахыг заахад ашиглагдана & Өөр сэдвийн төстэй хуудас дээгүүр гарч болно \\
\hline
Хүрээний нэмэлт оноо & Энгийн хайлт дээр нэмж асуулт аль хүрээнд хамаарахыг тодорхойлж, тэр хүрээний эх сурвалжид нэмэлт оноо өгнө. Бусдыг хасахгүй & Нэмэлт оноо хэт их бол буруу хүрээ давамгайлна \\
\hline
Салаагаар шүүх & Эхлээд тодорхойлсон салаандаа л хайна & Салааг буруу тодорхойлбол зөв эх сурвалж огт хайгдахгүй \\
\hline
\end{tabular}
\end{table}

Салаагаар шүүх хувилбарт салааг буруу тодорхойлбол зөв хариулт огт олдохгүй. Энэ ажлын мод 31 дэд сэдэвтэй, харин бодит асуултыг хараахан цуглуулаагүй тул нарийн ангилагч сургах боломжгүй. Иймд хүрээний нэмэлт онооны хувилбарыг сонгосон: асуултыг 11 том хүрээний түвшинд л ангилна, бусад эх сурвалжийг хасахгүй. Энэ нь модоор алхаж хайдаг арга биш, асуултыг ангилаад хайлтын оноонд зөөлөн нөлөөлдөг арга юм. Жишээ нь «элсэлтийн төлбөр хэд вэ» гэсэн асуултыг «Элсэлт» хүрээнд хамааруулж, элсэлтийн хуудсуудыг жагсаалтын дээд талд гаргана. Монгол хэл дээрх хайлтын нийтэд нээлттэй үнэлгээний багц 2026 оны 11-р сард хийсэн хайлтаар олдоогүй. Иймд энэ хувилбарыг ресепшнд асуудаг бодит асуултаар, нэмэлт оноогүй энгийн хайлтын тоотой хамт үнэлнэ (\ref{sec:exp-tree}-р хэсэг).
